\begin{theorem}[tangent cohomology and local miniversality]\hypertarget{tag-032U}{}\label{push:src:E.2.1}
\StatementRecord{push:src:E.2.1}{proved}{032U}{\SeriesRef{push:section:E.1}, \SeriesRef{push:src:E.3.1}}

For every sphere member with multiplicities $(3k,4l)$ and the
admissible logarithmic data of \ref{push:section:E.1}, put $z=kl$ and $d=3k+4l$.
On $\mathfrak B_{k,l}=\mathbb P(3k,4l)$ the root lift
$\bar f:X\to\mathfrak B_{k,l}$ satisfies
\[
 \begin{split}
 \mathcal E&=\mathcal O\oplus\mathcal O(-z),\\
 R^q\bar f_*\mathcal O_X&=
 \mathcal O,\quad\mathcal O\oplus\mathcal O(-z),\quad\mathcal O(-z),\\
 K_X&=\bar f^*\mathcal O(z-d).
 \end{split}
\]
The tangent cohomology is $(h^0,h^1,h^2,h^3)(X,T_X)=(1,1,1,0)$.
With the groupoid derivative convention
$\kappa(\Phi)=(\partial_c\Phi)\Phi^{-1}$, the class of
$\partial/\partial c$ is $e\otimes[\gamma]$.  The class of
$\partial/\partial u$ is
$-(2\pi i u)^{-1}e\otimes[\gamma]$.  Both are nonzero.
The Kuranishi germ is a smooth reduced line, and the fixed-data
period family is locally miniversal.
\end{theorem}


\begin{proof}
\begin{enumerate}[label={$\langle1\rangle$\arabic*.},leftmargin=*]
\item The root object is defined by the actual divisor identities
$f^*(0)=mS_1$ and $f^*(1)=nS_2$, with $m=3k$, $n=4l$.
It gives a holomorphic map to the analytic root stack.  Its base
changes to the two root discs are the actual smooth torus
families; the quotient actions are free.  At the cusp its charts
are $q=z_0z_1z_2$, so the source is smooth and the map is flat,
of relative dimension two, with Gorenstein fibres.  Properness is
checked on the root atlas, where it is properness of the compact
torus families and the original compact cusp filling.

\item Use weighted coordinates $x,y$ of weights $m,n$, with coarse
coordinate $t=y^m/(y^m-x^n)$.
The two root divisors have bundles $\mathcal O(n)$ and
$\mathcal O(m)$, and the cusp has bundle $\mathcal O(mn)$.
The root-forgetting morphism
\[
 r:\mathfrak B_{k,l}\longrightarrow\mathfrak B_{1,1},
 \qquad (x,y)\longmapsto(x^l,y^k),
\]
uses the group homomorphism $\lambda\mapsto\lambda^{kl}$ on
the weighted quotient.  It is a holomorphic morphism of analytic
stacks over $B$, with local maps $s\mapsto s^k$ and
$s\mapsto s^l$.  In particular $r^*\mathcal O(1)=\mathcal O(z)$.
For completeness, the analytic Picard group of this root curve has
the presentation
$\langle H,R_1,R_2\mid mR_1=H,\ nR_2=H\rangle$:
after subtracting the two stabilizer characters, a line bundle
descends to the coarse projective line.  Its only relations are
the two defining root identities.  The Smith form and
$\gcd(m,n)=1$ identify this group with $\mathbb Z$, with
$H=mn\mathcal O(1)$, $R_1=n\mathcal O(1)$, and
$R_2=m\mathcal O(1)$.  This proves the analytic torsion-free
Picard assertion used below.
Translations have identity vertical derivative and act trivially
on structure-sheaf cohomology.  Thus the vertical bundle computed
from the linear period factors is $r^*(\mathcal O\oplus
\mathcal O(-1))$.  The original determinant calculation can also
be repeated: its quotient line $L$ has
$L^{12}=\mathcal O(-p_0)=\mathcal O(-12z)$, so
$L=\mathcal O(-z)$ because the Picard group is torsion-free.
The fixed vector gives a subline $\mathcal O$; its extension
splits because $H^1(\mathcal O(z))=0$.

\item The cusp normalization resolution has terms
$\mathcal O_{\mathrm{dP}_6}$, three copies of
$\mathcal O_{\mathbb P^1}$, and two point sheaves.
Its global differential is $(a,b,c)\mapsto
(a-b+c,a-b+c)$, giving cohomology $(1,2,1)$; its alternating cup
pairing in degree one is nondegenerate.  The degree-one cocycle
$\gamma$ extends across the cusp because it vanishes on the toric
period lattice.  Its edge values $(0,-1,-1)$ have nonzero cup
product with $(1,0,-1)$.  On finite root charts its class is the
nonzero constant lattice cocycle $\gamma$, invariant under every
deck transformation.  These are explicit comparisons of the
same class, so it gives a nowhere-zero subline in
$G=R^1\bar f_*\mathcal O_X$.
The base-change and relative-duality argument \ref{push:src:E.3.1} applies on
these charts.  It gives $R^2\bar f_*\mathcal O_X=L$ and
$\det G=L$.  The extension of $L$ by the specified
$[\gamma]$-subline splits by the same $H^1(\mathcal O(z))=0$.
This proves all direct-image formulas.

\item The logarithmic tangent sequence is
\[
 0\longrightarrow\mathcal V=\bar f^*\mathcal E
 \longrightarrow\mathcal T=T_X(-\log W)
 \longrightarrow\mathcal H=\bar f^*\mathcal O(d-12z)
 \longrightarrow0.
\]
Its last map is the relative differential on smooth root charts
and sends $\sum a_i z_i\partial_{z_i}$ to
$(\sum a_i)q\partial_q$ at the cusp.  It is surjective there.
Since $K_{\mathfrak B}=\mathcal O(-d)$, determinants give
$K_X=\bar f^*\mathcal O(z-d)$.

\item We give the weighted cohomology calculation explicitly; simply
replacing $1$ by $kl$ in an acyclicity assertion would be incorrect.
The two quotient charts $x\ne0$, $y\ne0$ are finite quotients
of affine lines, and their overlap is a finite quotient of a
punctured line.  Finite averaging and Stein acyclicity compute
coherent cohomology by this two-chart complex.  Its homogeneous
Laurent monomials give
\[
 \begin{split}
 h^0\mathcal O(a)&=\#\{(i,j)\in\mathbb Z_{\ge0}^2:mi+nj=a\},\\
 h^1\mathcal O(a)&=h^0\mathcal O(-a-d).
 \end{split}
\]
The second equality identifies the remaining monomials
$x^{-i-1}y^{-j-1}$; it does not require a positivity theorem.
The bundles $\mathcal O(-z),\mathcal O(-2z)$ are acyclic.
Indeed a possible degree-one monomial would give
$m(i+1)+n(j+1)=az$, with $a=1$ or $2$.
Coprimality forces $i+1=lA$, $j+1=kB$, hence
$3A+4B=a$ with $A,B\ge1$, which is impossible.
Positive degree has no degree-one monomials, so
$H^1\mathcal O(z)=0$ as used above.

\item Put $b=d-12z$.  Both $b$ and $b-z$ are negative.
Degree-one monomials for $\mathcal O(b)$ would satisfy
$3A+4B=12$, with $A,B\ge1$, and none exist.
For $\mathcal O(b-z)$ the equation is $3A+4B=13$;
its only positive solution is $(A,B)=(3,1)$, producing
$x^{3l-2}y^{k-2}$ in the dual description.  It is allowed
exactly when $k>1$; write $e_k=1$ in that case and $e_k=0$
otherwise.  Projection formula and the two-column Leray sequence
now give
\[
 h^\bullet(\mathcal V)=(1,1,0,0),\qquad
 h^\bullet(\mathcal H)=(0,0,e_k,e_k),\qquad
 h^\bullet(\mathcal T)=(1,1,e_k,e_k).
\]
The program \texttt{e\_push\_exact.py} independently checks these
counts for all $1965$ coprime pairs with $k,l\le80$; the monomial
argument proves them for all pairs.

\item On the normalization $\nu:E=\mathrm{dP}_6\to W\subset X$,
the branchwise normal line is $K_E$: the invariant cusp form
$d\log x_1\wedge d\log x_2\wedge dq$ trivializes $K_X$
on the entire cusp neighbourhood.  Local normal coefficients
therefore give the exact sequence
\[
 0\longrightarrow\mathcal T\longrightarrow T_X
 \longrightarrow\nu_*K_E\longrightarrow0,
 \qquad h^\bullet(K_E)=(0,0,1).
\]
We must determine its connecting map when $k>1$.
The dual logarithmic sequence, tensored with $K_X$, has horizontal
subline $\bar f^*\mathcal O(13z-2d)$ and vertical quotient
$\bar f^*(\mathcal O(z-d)\oplus\mathcal O(2z-d))$.
The latter has no global sections, by the same equations
$3A+4B=1,2$ above.  Consequently
\[
 H^0(\Omega_X^1(\log W)\otimes K_X)
 =H^0(\mathfrak B,\mathcal O(13z-2d)).
\]
For $k>1$ this is generated by the monomial
$x^{3l-2}y^{k-2}$ multiplying the horizontal logarithmic
differential.  Both $x$ and $y$ are nonzero over $p_0$, so its
residue on $E$ is nonzero.  The ordinary logarithmic residue
sequence, tensored with $K_X$, identifies this residue map with
the Serre dual of
$H^2(E,K_E)\to H^3(X,\mathcal T)$.
One can check this identification directly in each chart
$q=z_0z_1z_2$: the residues of $dq/q$ are one on each branch,
and contraction with the invariant three-form is the normal
coefficient pairing.  Thus that connecting map has rank $e_k$.
The long exact sequence proves $h^2(T_X)=1$ and $h^3(T_X)=0$;
degrees zero and one are unchanged.

\item Finally $\partial_c\Pi=e\otimes\gamma$ gives the constant
groupoid cocycle $\lambda\mapsto\gamma(\lambda)e$ on every
regular torus.  On the enlarged finite roots the derivative of the
affine generator is $(r_j/m_j)e$, and the derivative of its
punctured connector is $(\log s_j)/(2\pi i)\,r_je$.
The identity $m_j(r_j/m_j)=\gamma(w_j)$ verifies the power
relation, while $\gamma A_j=\gamma$ verifies conjugation with
every lattice arrow.  Changing $\log s_j$ by $2\pi i N$
adds the derivative of the integral arrow $Nw_j$; the connector
formula is used only on the punctured overlap.  Its central
root-chart representative is the constant $(r_j/m_j)e$,
which is holomorphic.  At the cusp $\gamma$ kills the toric periods,
so the constant deck cocycle is exactly its extending cocycle
above.  Its image under the two isomorphisms
$H^1(\mathcal V)\simeq H^1(\mathcal T)\simeq H^1(T_X)$
is $e\otimes[\gamma]\ne0$.  The chain rule gives the asserted
$u$-class.  The Kuranishi local algebra is
$\mathbb C\{t\}/I$ since $h^1=1$.  The existing family induces
$t\mapsto g(c)$ with $g'(c_0)\ne0$.  The analytic inverse
function theorem makes this substitution injective; every
element of $I$ maps to zero, so $I=0$.  Its classifying map is
then an isomorphism of germs, proving local miniversality.
\end{enumerate}
\end{proof}
